% !TeX root = ../main.tex
% Add the above to each chapter to make compiling the PDF easier in some editors.

\chapter{Future Work}\label{chapter:future_work}

\section{Dedicated Benchmarking and Training Data}

The most significant limitation of this work is the training data.
We rely entirely on the ZeroShot benchmark data,
which was designed to enable cost estimation that generalizes to unseen database instances,
rather than to cover the full spectrum of PostgreSQL query behavior~\parencite{zeroshot}.

The ZeroShot training queries follow a narrow pattern of selective table scans,
equi-joins (joins on equality predicates, i.\,e., \texttt{JOIN ON a = b}), and a final aggregation to a single tuple~\parencite{zeroshot}.
In contrast, T3 generates training data from 16 distinct query structures,
explicitly including complex selection predicates
and combinations of filters, joins, aggregations, and sorts~\parencite{t3}.
These query types are largely absent from the ZeroShot training data,
which limits how well our model can generalize to more diverse workloads.

The JOB comparison was restricted to 77 of the 113 benchmark queries,
because the remaining queries exceed the 30\,s cutoff of the ZeroShot data~\parencite{zeroshot}.
Rieger and Neumann evaluated T3 on all 113 queries~\parencite{t3},
making a direct like-for-like comparison on the full JOB benchmark impossible.

Consequently, building a dedicated PostgreSQL benchmark in the style of T3's Umbra benchmark
would close the query diversity gap and enable a full JOB comparison.

The runtime distribution of the ZeroShot data is also considerably narrower than T3's Umbra training data~\parencite{t3}.
A dedicated benchmark covering a wider runtime range could allow the per-tuple prediction approach to be re-evaluated,
since T3's per-tuple normalization is specifically motivated by such a wide range~\parencite{t3}
and our ablation study could not test this under comparable conditions (Section~\ref{sec:evaluation_ablation}).

Moreover, a data gap is the absence of tablescan input cardinalities in the parsed ZeroShot plans.
Our implementation always sets this value to 1, and the model compensates via the many other features.
The estimated row count per table is available in PostgreSQL's system catalog \parencite{postgresql_docs_pg_class}.
Future work should extend the data collection to record this information.

\section{Feature Vector}

\paragraph{Finer-Grained Pipeline Features.}
Currently, the pipeline-level block of our feature vector aggregates features across all nodes in a pipeline.
For example, the planned workers are summed over the pipeline.
A natural extension is to expand the pipeline block with per-node entries,
at the cost of a larger feature vector.
For example, this would let the model clearly distinguish a pipeline where only one node is parallelized
from one where all nodes run in parallel.

\paragraph{Purely Query-Level Feature Vector.}
Per-query prediction outperforms per-pipeline and per-tuple prediction in our setting.
Therefore, a future direction could be to design the feature vector directly at the query level,
without the intermediate pipeline decomposition.
This would further sidestep the PostgreSQL pipeline timing problem identified in the ablation study
(Section~\ref{sec:evaluation_ablation}), while also simplifying the model pipeline.

\paragraph{Joint Pipeline and Query Context for Pipeline-Level Prediction.}
For the per-pipeline prediction approach, a possible extension is to provide
two feature vectors per pipeline: one capturing only the local pipeline features,
and one summarizing the full query context.
This gives the model information about how a pipeline relates to the rest of the query,
which may improve per-pipeline runtime prediction.

\paragraph{PostgreSQL-Specific Features.}
We include PostgreSQL planner cost estimates as features.
However, the ablation study confirms that the PostgreSQL cost estimator is highly inaccurate
for complex queries (median Q-Error above 1000, see Table~\ref{tab:median_qerror_job}).
A targeted ablation study that removes or replaces these planner-cost features
could determine whether they help or hurt prediction accuracy.
More broadly, to assess the importance of each PostgreSQL-native feature,
future work can use LightGBM's built-in feature importance scores.
Thereby, it could identify which features drive the model's predictions the most.


\paragraph{Mitigating the Pipeline Timing Issue.}
The pipeline timing problem identified in the ablation study arises because
PostgreSQL's \texttt{EXPLAIN ANALYZE} reports cumulative wall-clock intervals per node,
not exclusive execution times~\parencite{postgresql_docs_explain}.
A node's reported interval includes time spent waiting for its children,
which causes intervals to overlap across pipeline boundaries.

Hence, a thorough long-term approach would be to modify PostgreSQL's executor instrumentation to track exclusive CPU time per node directly.
PostgreSQL already exposes per-node timing hooks through its instrumentation layer~\parencite{postgresql_docs_explain}.
Extending these hooks to subtract child contribution at collection time would yield exclusive per-node times.
Such instrumentation could then serve as the basis for a dedicated benchmark that collects richer timing data than the current ZeroShot dataset provides.

\section{Prediction Latency}

Since both our model and T3 use LightGBM gradient-boosted decision trees,
prediction should be fast in practice.
Rieger and Neumann note that T3's low prediction latency is a key practical advantage~\parencite{t3}.
Future work should benchmark the prediction latency of our PostgreSQL model
to confirm its performance,
and to compare it directly against T3's reported latency~\parencite{t3}.
A concrete step toward this would be to compile our trained LightGBM model to native machine code using the \texttt{lleaves} framework, as T3 does~\parencite{t3}.
T3 already demonstrates that this compilation step reduces inference latency from 22\,\textmu{}s to approximately 4\,\textmu{}s, and our model could benefit from the same approach.

% --- ORIGINAL NOTES (for review) ---
%
% -Dedicated Benchmarking with Queries from T3 paper for PG, we had time limitations:
%
% Why?
%
% Data from Zeroshot kann limitierender faktor sein.
% query distribution vielleicht nicht so gut wie bei t3 umbra, zu wenig categories, auch timewise nicht so stark distributed -- check und vergleiche genau, inwieweit daten überlegen sind.
% INPUT CARD existiert für scans nicht in zeroshot data, dadruch immer auf 1 gelegt.
%
% JOB nur 77 von 113 queries in zeroshot, comparison eingeschränkt.
%
%
% - Änderungen am Feature Vector:
%
%
% Pipeline-level block aufbrechen, und grössere feature size in kauf nehmen. z.b. planned parallelism per node, etc...
%
% Weil query-level am besten funktioniert hat: Versuchen auch mal weg von pipeline level zu gehen und feature vector auf per query level zu machen
%
% -> dadurch könnte man weiterhin PG Problem umgehen, dass pipeline prediction targets aufgrund der timing issues (siehe eval) nicht präzise genug ist.
% Oder man kann sich natürlich anschauen, wie man dieses pipeline time issue mitigaten kann.
%
% Interessante future work könnte auch sein, dem Model zwei feature vectors für die per pipeline prediction zu geben. Einmal nur die features per pipeline und einmal alle query features.
%
% Impact der inkludierten pg Cost estimation im vektor untersuchen, haben in eval gesehen, dass der PG cost estimator noch viele Schwierigkeiten hat
%
%
% - Untersuche Prediction Speed/latency, dadurch, dass wir LightGBM trees benutzt haben, sollte es wie in Umbra T3, schnell sein
